%% Introduction generale — 1 page
\chapter*{Introduction générale}
\addcontentsline{toc}{chapter}{Introduction générale}

L'électrification accélérée du parc automobile mondial place le \emph{Battery
Management System} au c\oe{}ur des défis ingénieurs des prochaines décennies.
Valider un \bms couvre aujourd'hui des centaines d'exigences réparties
entre le modèle de cellule, les algorithmes embarqués, la génération de
code et les tests d'intégration. Reliés par un V-model
rigoureux, ces activités sont encore trop souvent menées manuellement,
avec des outils hétérogènes, sans traçabilité automatisée et sans
capacité de synthèse intelligente.

Ce projet de fin d'études propose et implémente un \textbf{outil intelligent
d'assistance à la conception et à la validation des \bms}, fondé sur un
V-model dual~: une branche Modèle produit et teste un \bms complet
(\emph{plant} ECM~2RC, maître, esclaves, prédicteur LSTM) sous
\textsc{Simulink}~; une branche Logicielle produit un pipeline
automatisé (Validateur, Analyseur, Générateur de rapports LLM,
Orchestrateur, Interface utilisateur). La convergence des deux branches
constitue la démonstration finale de la valeur ajoutée de l'approche.

Le rapport est structuré en cinq chapitres. Le \textbf{chapitre~1}
pose le cadre industriel et l'état de l'art (V-model, outils de
validation, apports de l'IA). Le \textbf{chapitre~2} analyse les
besoins et spécifie les 86~exigences du système. Le \textbf{chapitre~3}
décrit la branche Modèle et présente ses résultats de validation
(MIL, SIL, PIL). Le \textbf{chapitre~4} décrit la branche Logicielle
et ses 45~tests unitaires. Le \textbf{chapitre~5} documente
l'intégration des deux branches et la validation d'acceptation
complète.
